% =====================================================
% settings.tex
% Cấu hình chung cho báo cáo
% =====================================================

% =========================
% 1. CẤU HÌNH TIẾNG VIỆT VÀ FONT TIMES NEW ROMAN TOÀN BÀI
% =========================
\usepackage{fontspec}
\IfFileExists{fonts/times.ttf}{%
  \setmainfont{times.ttf}[
    Path = fonts/,
    BoldFont = timesbd.ttf,
    ItalicFont = timesi.ttf,
    BoldItalicFont = timesbi.ttf
  ]%
  \setsansfont{times.ttf}[
    Path = fonts/,
    BoldFont = timesbd.ttf,
    ItalicFont = timesi.ttf,
    BoldItalicFont = timesbi.ttf
  ]%
  \setmonofont{times.ttf}[
    Path = fonts/,
    BoldFont = timesbd.ttf,
    ItalicFont = timesi.ttf,
    BoldItalicFont = timesbi.ttf
  ]%
}{%
  \IfFontExistsTF{Times New Roman}{%
    \setmainfont{Times New Roman}%
    \setsansfont{Times New Roman}%
    \setmonofont{Times New Roman}%
  }{%
    \setmainfont{Tinos}%
    \setsansfont{Tinos}%
    \setmonofont{Tinos}%
  }%
}

\usepackage{polyglossia}
\setdefaultlanguage{vietnamese}

% Tự động ngắt dòng mềm mại, chống tràn lề phải và chống kéo giãn khoảng trắng
\tolerance=2000
\emergencystretch=4em
\hbadness=10000
\raggedbottom

% =========================
% 2. CẤU HÌNH TRANG, CỠ CHỮ 13PT VÀ GIÃN DÒNG 1.5
% =========================
\usepackage[
    a4paper,
    left=3.0cm,
    right=1.5cm,
    top=2.0cm,
    bottom=2.0cm
]{geometry}

\usepackage{setspace}
\onehalfspacing

% Cỡ chữ mặc định của toàn bộ nội dung văn bản là 13pt (baseline chuẩn 15.6pt,
% khi nhân với \onehalfspacing sẽ đúng giãn dòng 1.5 mà không làm thưa bảng).
\makeatletter
\renewcommand{\normalsize}{%
    \@setfontsize\normalsize{13pt}{15.6pt}%
    \abovedisplayskip 10pt plus 2pt minus 5pt%
    \abovedisplayshortskip 0pt plus 3pt%
    \belowdisplayshortskip 6pt plus 3pt minus 3pt%
    \belowdisplayskip \abovedisplayskip
    \let\@listi\@listI
}
\renewcommand{\small}{\@setfontsize\small{13pt}{15.6pt}}
\renewcommand{\footnotesize}{\@setfontsize\footnotesize{13pt}{15.6pt}}
\makeatother
\normalsize

% Không thụt lề đoạn văn, thẳng bằng lề trái 100%
\setlength{\parindent}{0pt}
\setlength{\parskip}{5pt}

% =========================
% 3. GÓI HỖ TRỢ, CỘT BẢNG CHUẨN & DANH SÁCH THẲNG LỀ TRÁI
% =========================
\usepackage{graphicx}
\usepackage{float}
\usepackage{array}
\usepackage{longtable}
\usepackage{booktabs}
\usepackage{multirow}
\usepackage{tabularx}
\usepackage{caption}
\usepackage{amsmath, amssymb}
\usepackage{enumitem}
\usepackage{etoolbox}

% Định nghĩa các kiểu cột bảng căn trái (L, Y) và căn giữa (C) để KHÔNG BAO GIỜ
% bị kéo giãn khoảng cách chữ quá xa trong các ô bảng hẹp:
\newcolumntype{L}[1]{>{\raggedright\arraybackslash}p{#1}}
\newcolumntype{C}[1]{>{\centering\arraybackslash}p{#1}}
\newcolumntype{R}[1]{>{\raggedleft\arraybackslash}p{#1}}
\newcolumntype{Y}{>{\raggedright\arraybackslash}X}

% Trong bảng (table/tabularx), tự động dùng giãn dòng đơn chuẩn (1.2) để các dòng
% trong bảng gọn gàng, sát nhau vừa phải, không bị cách xa nhau:
\AtBeginEnvironment{table}{\singlespacing\renewcommand{\arraystretch}{1.2}}
\AtBeginEnvironment{tabular}{\singlespacing\renewcommand{\arraystretch}{1.2}}
\AtBeginEnvironment{tabularx}{\singlespacing\renewcommand{\arraystretch}{1.2}}
\AtBeginEnvironment{longtable}{\singlespacing\renewcommand{\arraystretch}{1.2}}

% Tất cả các gạch đầu dòng và đánh số đều sát bằng lề trái (leftmargin=*), không thụt vào
\setlist[itemize]{
    label=-,
    leftmargin=*,
    align=left,
    labelsep=0.5em,
    topsep=2pt,
    itemsep=2pt,
    parsep=0pt
}
\setlist[enumerate]{
    leftmargin=*,
    align=left,
    labelsep=0.5em,
    topsep=2pt,
    itemsep=2pt,
    parsep=0pt
}

\usepackage{titlesec}
\usepackage{tocloft}
\usepackage{hyperref}
\usepackage{tikz}
\usetikzlibrary{shapes.geometric, arrows.meta, positioning}
\usepackage{listings}
\usepackage{xcolor}
\usepackage{eso-pic}
\usepackage{fancyhdr}

\lstset{
    language=Python,
    basicstyle=\rmfamily\fontsize{13pt}{16pt}\selectfont,
    keywordstyle=\color{blue!80!black}\bfseries,
    commentstyle=\color{green!50!black}\itshape,
    stringstyle=\color{red!70!black},
    numbers=left,
    numberstyle=\fontsize{11pt}{13pt}\selectfont\color{gray},
    xleftmargin=1.8em,
    framexleftmargin=1.5em,
    frame=single,
    breaklines=true,
    breakatwhitespace=false,
    columns=flexible,
    keepspaces=true,
    showstringspaces=false,
    tabsize=4
}

% =========================
% 4. HYPERLINK
% =========================
\hypersetup{
    colorlinks=true,
    linkcolor=black,
    citecolor=black,
    urlcolor=blue
}

% Số trang đặt giữa phía dưới (footer); áp dụng cả cho trang mở đầu chương.
\setlength{\headheight}{17pt}
\setlength{\footskip}{24pt}
\pagestyle{fancy}
\fancyhf{}
\fancyfoot[C]{\fontsize{13pt}{15pt}\selectfont\thepage}
\renewcommand{\headrulewidth}{0pt}
\renewcommand{\footrulewidth}{0pt}
\fancypagestyle{plain}{
    \fancyhf{}
    \fancyfoot[C]{\fontsize{13pt}{15pt}\selectfont\thepage}
    \renewcommand{\headrulewidth}{0pt}
    \renewcommand{\footrulewidth}{0pt}
}

% =========================
% 5. ĐÁNH SỐ HÌNH, BẢNG, SƠ ĐỒ THEO CHƯƠNG & CHÚ THÍCH CỠ 11PT
% =========================
\counterwithin{table}{chapter}
\counterwithin{figure}{chapter}
\renewcommand{\thetable}{\thechapter.\arabic{table}}
\renewcommand{\thefigure}{\thechapter.\arabic{figure}}

% Môi trường riêng cho sơ đồ để đánh số và lập danh mục tự động độc lập.
\newfloat{sodo}{H}{los}[chapter]
\floatname{sodo}{Sơ đồ}
\renewcommand{\thesodo}{\thechapter.\arabic{sodo}}
\makeatletter
\newcommand{\l@sodo}{\@dottedtocline{1}{0pt}{2.8cm}}
\makeatother

% Chú thích ảnh, bảng, sơ đồ để cỡ chữ 11pt theo đúng yêu cầu
\DeclareCaptionFont{caption11pt}{\fontsize{11pt}{13.5pt}\selectfont}
\captionsetup{
    font=caption11pt,
    labelfont=bf,
    textfont=normalfont,
    labelsep=period,
    justification=centering,
    aboveskip=6pt,
    belowskip=6pt
}
\captionsetup[table]{name=Bảng,position=top}
\captionsetup[figure]{name=Hình,position=bottom}
\captionsetup[sodo]{name=Sơ đồ,position=bottom}

\renewcommand{\listtablename}{Danh mục bảng}
\renewcommand{\listfigurename}{Danh mục hình}
\renewcommand{\contentsname}{MỤC LỤC}
\renewcommand{\bibname}{TÀI LIỆU THAM KHẢO}

% =========================
% 6. ĐỊNH DẠNG TIÊU ĐỀ CHƯƠNG (CỠ 16PT, IN HOA) & CÁC MỤC CON (13PT)
% =========================
\titleformat{\chapter}
  [display]
  {\centering\bfseries\fontsize{16pt}{21pt}\selectfont}
  {CHƯƠNG \thechapter}
  {0.25cm}
  {\MakeUppercase}

\titlespacing*{\chapter}{0pt}{0pt}{16pt}

\titleformat{\section}
  {\bfseries\fontsize{13pt}{16pt}\selectfont}
  {\thesection.}
  {0.5em}
  {}
\titlespacing*{\section}{0pt}{10pt}{4pt}

\titleformat{\subsection}
  {\bfseries\itshape\fontsize{13pt}{16pt}\selectfont}
  {\thesubsection.}
  {0.5em}
  {}
\titlespacing*{\subsection}{0pt}{8pt}{3pt}

\titleformat{\subsubsection}
  {\itshape\fontsize{13pt}{16pt}\selectfont}
  {\thesubsubsection.}
  {0.5em}
  {}
\titlespacing*{\subsubsection}{0pt}{6pt}{2pt}

\setcounter{secnumdepth}{3}
\setcounter{tocdepth}{3}

% =========================
% 7. ĐỊNH DẠNG MỤC LỤC & DANH MỤC TỰ ĐỘNG (DẤU CHẤM DẪN TRANG CHUẨN)
% =========================
\renewcommand{\cftchapfont}{\bfseries\fontsize{13pt}{16pt}\selectfont}
\renewcommand{\cftchappagefont}{\bfseries\fontsize{13pt}{16pt}\selectfont}
\renewcommand{\cftchappresnum}{CHƯƠNG~}
\renewcommand{\cftchapaftersnum}{.}
\setlength{\cftchapnumwidth}{2.8cm}

\renewcommand{\cftsecfont}{\bfseries\fontsize{13pt}{16pt}\selectfont}
\renewcommand{\cftsecpagefont}{\bfseries\fontsize{13pt}{16pt}\selectfont}
\renewcommand{\cftsecaftersnum}{.}

\renewcommand{\cftsubsecfont}{\itshape\fontsize{13pt}{16pt}\selectfont}
\renewcommand{\cftsubsecpagefont}{\fontsize{13pt}{16pt}\selectfont}
\renewcommand{\cftsubsecaftersnum}{.}

\renewcommand{\cftsubsubsecfont}{\fontsize{13pt}{16pt}\selectfont}
\renewcommand{\cftsubsubsecpagefont}{\fontsize{13pt}{16pt}\selectfont}
\renewcommand{\cftsubsubsecaftersnum}{.}

\renewcommand{\cftchapleader}{\cftdotfill{\cftdotsep}}
\setlength{\cftbeforechapskip}{5pt}
\setlength{\cftsecindent}{0pt}
\setlength{\cftsubsecindent}{0pt}
\setlength{\cftsubsubsecindent}{0pt}

\renewcommand{\cfttoctitlefont}{\hfill\bfseries\fontsize{16pt}{21pt}\selectfont\MakeUppercase}
\renewcommand{\cftaftertoctitle}{\hfill}
\setlength{\cftbeforetoctitleskip}{0pt}
\setlength{\cftaftertoctitleskip}{16pt}

% Cấu hình mục lục tự động cho Bảng, Sơ đồ và Hình (có dấu chấm dẫn số trang tự động)
\renewcommand{\cfttabfont}{\fontsize{13pt}{16pt}\selectfont}
\renewcommand{\cfttabpagefont}{\fontsize{13pt}{16pt}\selectfont}
\renewcommand{\cfttabpresnum}{Bảng~}
\renewcommand{\cfttabaftersnum}{.}
\setlength{\cfttabnumwidth}{2.2cm}
\setlength{\cfttabindent}{0pt}

\renewcommand{\cftfigfont}{\fontsize{13pt}{16pt}\selectfont}
\renewcommand{\cftfigpagefont}{\fontsize{13pt}{16pt}\selectfont}
\renewcommand{\cftfigpresnum}{Hình~}
\renewcommand{\cftfigaftersnum}{.}
\setlength{\cftfignumwidth}{2.2cm}
\setlength{\cftfigindent}{0pt}

\makeatletter
\newcommand{\AutoListSection}[3]{%
    {\bfseries\fontsize{13pt}{16pt}\selectfont #1\par}
    \vspace{4pt}
    \begingroup
    \fontsize{13pt}{16pt}\selectfont
    \def\numberline##1{#2~##1.\enspace}%
    \@starttoc{#3}%
    \endgroup
}

\newcommand{\CombinedLists}{%
    \chapter*{DANH MỤC CÁC BẢNG, SƠ ĐỒ VÀ HÌNH}
    \addcontentsline{toc}{chapter}{DANH MỤC CÁC BẢNG, SƠ ĐỒ VÀ HÌNH}
    \AutoListSection{1. Danh mục bảng}{Bảng}{lot}
    \vspace{12pt}
    \AutoListSection{2. Danh mục sơ đồ}{Sơ đồ}{los}
    \vspace{12pt}
    \AutoListSection{3. Danh mục hình}{Hình}{lof}
}
\makeatother

% =========================
% 8. THÔNG TIN SINH VIÊN/BÁO CÁO
% =========================
\newcommand{\BoGiaoDuc}{BỘ GIÁO DỤC VÀ ĐÀO TẠO}
\newcommand{\TenTruong}{TRƯỜNG ĐẠI HỌC ĐẠI NAM}

\newcommand{\TenSinhVien}{Nhóm 11 (Đề tài 13)}
\newcommand{\MaSinhVien}{FIT4013 -- Nhóm 11}
\newcommand{\KhoaHoc}{2026 -- 2027}

\newcommand{\TenBaoCao}{ĐỀ TÀI 13: DỰ ĐOÁN NHÂN VIÊN CÓ KHẢ NĂNG NGHỈ VIỆC\\(HR EMPLOYEE ATTRITION)}
\newcommand{\LoaiBaoCao}{BÁO CÁO BÀI TẬP LỚN}
\newcommand{\HocPhan}{HỌC PHẦN NHẬP MÔN HỌC MÁY}
\newcommand{\Nganh}{CÔNG NGHỆ THÔNG TIN}
\newcommand{\ChuyenNganh}{KHOA HỌC MÁY TÍNH / HỌC MÁY (FIT4013)}

\newcommand{\GiangVienHD}{ThS. Lê Thị Thùy Trang}

\newcommand{\DiaDiem}{HÀ NỘI}
\newcommand{\Nam}{2026}

% Logo trường. Có thể thay bằng file khác trong thư mục images/.
\newcommand{\LogoTruong}{images/logo-dai-nam.png}

% =========================
% 9. LỆNH CHƯƠNG KHÔNG ĐÁNH SỐ (CỠ 16PT, IN HOA) & NGUỒN CHÚ THÍCH CỠ 11PT
% =========================
\newcommand{\UnnumberedChapter}[1]{
    \chapter*{\MakeUppercase{#1}}
    \addcontentsline{toc}{chapter}{\MakeUppercase{#1}}
}

% Ghi nguồn ngay dưới bảng, hình hoặc sơ đồ với cỡ chữ 11pt đồng bộ chú thích
\newcommand{\Nguon}[1]{%
    \par\vspace{3pt}
    {\centering\fontsize{11pt}{13.5pt}\selectfont\itshape Nguồn: #1\par}
}

% Với môi trường sodo của gói float, caption luôn được dồn xuống cuối (cỡ 11pt).
\newcommand{\SodoCaption}[2]{%
    \caption[#1]{#1\\[2pt]{\fontsize{11pt}{13.5pt}\selectfont\textit{Nguồn: #2}}}%
}

% =========================
% 10. KHUNG TRANG BÌA
% =========================
\newcommand{\DrawCoverBorder}{
\AddToShipoutPictureBG*{
\begin{tikzpicture}[remember picture, overlay]
    \draw[line width=1pt]
        ([xshift=1.2cm,yshift=-1.2cm]current page.north west)
        rectangle
        ([xshift=-1.2cm,yshift=1.2cm]current page.south east);
\end{tikzpicture}
}
}
